%==============================================================================
\section{Computational study}\label{sec:comp}
%==============================================================================

Computation here is how the theory is tested. The quantity the
structural results single out, the coverage bound $q$, is also the root relaxation of
every configuration-based method, so a benchmark that records optima and bounds
measures the theory's central variable directly. The campaign also settles a question
that predates the theory and motivated the Benders solver in the first place: the SSP
has a polynomially solvable loading subproblem, which is the textbook precondition for
Benders decomposition, and whether a decomposition can exploit it against modern
compact models had not been tested.

\subsection{Experimental design}\label{ssec:bench}

\paragraph{Instances.}
The benchmark uses the standard SSP instance families, taken from the collection
assembled by \citet{Otiai2024}: the sets of \citet{Catanzaro2015}, of
\citet{Crama1994}, and four of the sets of \citet{Laporte2004}. Every instance in
every family is run; none is skipped, and no instance is selected on the basis of an
outcome. Table~\ref{tab:families} gives the sizes.

\begin{table}[h]\centering
\footnotesize
\begin{tabular}{@{}lrrrl@{}}
\toprule
Family & Instances & Jobs $n$ & Tools $\lvert T\rvert$ & Source \\
\midrule
Catanzaro & $171$ & $8$--$40$  & $5$--$60$  & \citet{Catanzaro2015} \\
Crama     & $160$ & $10$--$40$ & $10$--$60$ & \citet{Crama1994} \\
Laporte3  & $340$ & $8$        & $15$--$25$ & \citet{Laporte2004} \\
Laporte4  & $330$ & $9$        & $15$--$25$ & \citet{Laporte2004} \\
Laporte5  & $340$ & $15$       & $15$--$25$ & \citet{Laporte2004} \\
Laporte7  & $80$  & $10$--$15$ & $10$--$20$ & \citet{Laporte2004} \\
\midrule
Total     & $1{,}421$ & & & \\
\bottomrule
\end{tabular}
\caption{The instance families. An instance is a matrix together with a magazine
capacity. The Crama collection publishes the same matrices at four different
capacities, so its $40$ matrices give $160$ instances; the same convention is applied
throughout.}
\label{tab:families}
\end{table}

\paragraph{Protocol.}
Every run of the Benders solver and of the compact baselines is given a single uniform
time limit of one hour and $16$ gigabytes of memory. There is no early stopping and no
per-configuration retirement rule. All other solver parameters are left at their
defaults, so that the comparison measures the formulations rather than parameter
tuning. The branch-and-price prototypes are run separately, under a ten-minute limit,
for the reason given in Section~\ref{ssec:results-bnp}.

\paragraph{Hardware and software.}
All runs were executed on the LIMOS SLURM cluster, confined to a homogeneous pool of
identical nodes (AMD EPYC~7452, $2\times32$ cores, $512$\,GB of memory), so that
wall-clock times are comparable across runs. The mixed-integer solver is CPLEX~22.1.1
throughout, pinned to four dedicated threads per run. The branch-and-price prototypes
run single-threaded under PySCIPOpt. The campaign was distributed over $120$ SLURM
array shards and produced $16{,}994$ individual runs.

\paragraph{Configurations.}
Twelve solver configurations were run. Three are the reimplemented compact baselines:
F4 of \citet{Catanzaro2015}, written \texttt{CATZ-F4}; LSS of \citet{Laporte2004}; and
SSPMF of \citet{daSilva2024}. The remaining nine are the Benders solver, in an ablation
over its cut families and its strengthenings. The naming is as follows:
\texttt{BBC-LP} uses the dual optimality cuts, \texttt{BBC-K} the combinatorial cuts,
\texttt{+T} adds the triplet bounds, \texttt{+F} enables fractional separation,
\texttt{+C} the conflict-graph row, \texttt{+H} the hybrid-genetic incumbent,
\texttt{+P} the Pareto-optimal cut lifting, and \texttt{+ACC} all three strengthenings
together.

\paragraph{The reported quantity.}
The common metric is the empty-start cost of the sequence a method returns, evaluated
by the loading oracle. Each formulation has its own native objective, and two of them
use the free-initial convention, so the native values are not comparable across
methods. The empty-start cost of the returned sequence is defined for every method
regardless of its formulation, and it is what makes the cross-method comparison of
Section~\ref{ssec:validation} meaningful.

\subsection{Validation}\label{ssec:validation}

Correctness was established before any result was recorded, and again from the results
themselves.

\paragraph{Before the campaign.}
Every exact solver was checked against brute-force optima on a bank of small instances,
in every combination of its option flags. The loading oracle itself was verified
against an exact dynamic program over magazine states; the verification described in
Appendix~\ref{app:verify} repeats this on $83{,}670$ instance-and-sequence pairs and
finds no disagreement.

\paragraph{Identifying an instance.}
An instance is identified by the family it comes from, its file name, its number of
jobs, its number of tools, and its magazine capacity. All five are needed. The Crama
collection reuses the same file names at four different capacities, so two runs that
agree on the name can be runs on different instances. An earlier version of this
analysis keyed on the name alone; it reported disagreements between methods that, on
inspection, were comparisons between an instance at one capacity and the same matrix at
another. The comparison below uses the full identifier.

\paragraph{Cross-method agreement.}
Of the $1{,}421$ instances, $1{,}115$ were solved to proven optimality by at least one
method and $999$ by at least two. On those $999$ there are $41{,}673$ pairs of methods
that both proved an optimum. \emph{In none of them do the two disagree.} Every pair of
methods that both closed an instance returned a sequence of the same empty-start cost.

\paragraph{The two conventions, measured.}
Proposition~\ref{prop:conv} states that the empty-start and free-initial costs of a
sequence differ by the size of the initial load. The campaign confirms it on real
instances. The SSPMF model reports the free-initial cost natively, and on all
$1{,}038$ instances it closed, the difference between its reported value and the
empty-start cost of the same sequence equals the magazine capacity $b$ on $1{,}029$ of
them. On the other nine it is smaller than $b$, and those are exactly the instances on
which an optimal sequence begins with a job that needs fewer than $b$ tools, so the
magazine is not filled at the start. The two baselines that report the empty-start cost
natively, \texttt{CATZ-F4} and LSS, show a difference of zero on all $1{,}729$ runs in
which they closed an instance.

\paragraph{Tools that no job needs.}
Four matrices in the collection contain a tool that appears in no job's requirement
set: Crama \texttt{s2n007} and \texttt{s2n009}, Laporte3 \texttt{L28-7} and Laporte5
\texttt{L6-8}. Counted with the capacities at which they appear, that is ten of the
$1{,}421$ instances. The coverage bound of Corollary~\ref{cor:lb} must therefore be read as
$q=\lvert\U\rvert-b$, using the set $\U$ of tools that some job actually requires, and
not as $\lvert T\rvert-b$ using the tool count declared in the instance file. On those
ten the two differ, and the smaller value is the correct one: on Laporte3
instance \texttt{L28-7}, which declares $25$ tools of which one is never required and
has $b=5$, the optimum in the free-initial convention is $19$, which is
$\lvert\U\rvert-b=24-5$ and not $\lvert T\rvert-b=20$. All coverage-bound figures in
this section use $\lvert\U\rvert$.

\paragraph{A defect found and corrected.}
The first campaign gave the Benders solver poor results. The cause was that its master
problem lacked the coverage row $\theta\ge\lvert\U\rvert$ of Section~\ref{ssec:bbc},
which every other method carries implicitly through its formulation. The row was added
and the whole campaign re-run. All numbers below come from the corrected solver. The
episode is reported because it is the reason the earlier interim results were
superseded, and because it illustrates the general point of Section~\ref{sec:disc}: on
this problem a missing bound is expensive in a way that a missing search improvement is
not.

\paragraph{Runs that failed.}
Of the $16{,}994$ runs, $427$ returned no result. Three hundred and seventy-eight of
these are crashes of the worker subprocess and $49$ are a failure of the harness itself
when closing its log file on the cluster filesystem. The failures are not spread
evenly: $352$ of the $378$ crashes belong to the LSS baseline, and $239$ of those to
the Laporte5 family alone, whose $15$-job instances give LSS its largest models. LSS is
therefore reported against a denominator of $1{,}069$ rather than $1{,}421$, and its
success rate is not comparable with the others, because the runs it lost are
concentrated on the family it finds hardest. The remaining $75$ failures are scattered
across the nine Benders configurations, at most $26$ in any one of them.

\subsection{How far each method gets}\label{ssec:results-solved}

Table~\ref{tab:solves} reports, for each configuration and each family, the number of
instances solved to proven optimality against the number of runs that returned a
result.

\begin{table}[t]\centering
\scriptsize
\setlength{\tabcolsep}{4pt}
\begin{tabular}{@{}lrrrrrrr@{}}
\toprule
Configuration & Catanzaro & Crama & Laporte3 & Laporte4 & Laporte5 & Laporte7 & All\\
\midrule
SSPMF              & $73/171$ & $68/160$ & $340/340$ & $330/330$ & $161/336$ & $66/80$ & $1038/1417$\\
\texttt{CATZ-F4}   & $55/171$ & $44/160$ & $340/340$ & $330/330$ & $72/327$  & $44/78$ & $885/1406$\\
LSS                & $52/132$ & $42/124$ & $340/340$ & $330/330$ & $39/101$  & $41/42$ & $844/1069$\\
\midrule
\texttt{BBC-LP+T}  & $38/168$ & $34/156$ & $340/340$ & $192/330$ & $84/336$  & $39/80$ & $727/1410$\\
\texttt{BBC-LP}    & $38/169$ & $35/151$ & $340/340$ & $191/329$ & $83/338$  & $39/80$ & $726/1407$\\
\texttt{BBC-K}     & $39/166$ & $34/145$ & $340/340$ & $190/329$ & $83/335$  & $39/80$ & $725/1395$\\
\texttt{BBC-LP+F+H}& $43/171$ & $37/160$ & $316/340$ & $157/330$ & $97/339$  & $40/80$ & $690/1420$\\
\texttt{BBC-LP+F}  & $37/169$ & $32/160$ & $313/340$ & $156/329$ & $85/340$  & $36/80$ & $659/1418$\\
\texttt{BBC-LP+F+C}& $37/170$ & $32/160$ & $311/340$ & $157/329$ & $84/336$  & $36/79$ & $657/1414$\\
\texttt{BBC-LP+ACC}& $43/170$ & $37/160$ & $286/340$ & $154/330$ & $97/339$  & $39/80$ & $656/1419$\\
\texttt{BBC-K+F}   & $37/171$ & $21/121$ & $316/340$ & $156/329$ & $86/338$  & $35/79$ & $651/1378$\\
\texttt{BBC-LP+F+P}& $35/171$ & $32/160$ & $282/340$ & $155/328$ & $83/336$  & $35/79$ & $622/1414$\\
\bottomrule
\end{tabular}
\caption{Instances solved to proven optimality, against runs that returned a result,
under a one-hour limit. Denominators below the family size are the failed runs of
Section~\ref{ssec:validation}; the LSS row is the one where this matters, and its
totals should not be read against the others.}
\label{tab:solves}
\end{table}

Three things can be read from Table~\ref{tab:solves}.

The multicommodity model of \citet{daSilva2024} is the strongest single method, as its
authors report. That their claim replicates on an independent implementation and a
different instance collection is worth recording in its own right. It closes $1{,}038$
instances; the best Benders configuration closes $727$.

The families separate the methods rather than rank them uniformly. Laporte3 is solved
completely by six of the twelve configurations. Laporte4 splits them sharply: the
three compact models solve all $330$, and the Benders solver solves between $154$ and
$192$. Catanzaro and Crama are hard for everything, and no method closes more than
half of either.

Among the Benders variants, the three that do not separate at fractional points
(\texttt{BBC-LP}, \texttt{BBC-LP+T}, \texttt{BBC-K}) are the three best, within two
instances of one another. Every configuration carrying \texttt{+F} is worse than all
three. Section~\ref{ssec:results-bbc} takes that apart.

\paragraph{Solving times.}
Table~\ref{tab:times} reports times on the $535$ instances that all eight of the listed
configurations solve. The shifted geometric mean with a shift of ten seconds is used,
following standard practice for exact-method comparisons, because it is insensitive to
both very small and very large individual values.

\begin{table}[h]\centering
\footnotesize
\begin{tabular}{@{}lrrr@{}}
\toprule
Method & shifted geometric mean (s) & median (s) & maximum (s)\\
\midrule
SSPMF              & $1.90$  & $0.140$  & $48.7$\\
\texttt{BBC-K}     & $14.50$ & $0.018$  & $845.8$\\
\texttt{BBC-LP+T}  & $16.03$ & $0.066$  & $1126.0$\\
\texttt{BBC-LP}    & $16.11$ & $0.067$  & $1128.4$\\
\texttt{BBC-LP+F+H}& $30.74$ & $0.003$  & $3584.6$\\
\texttt{BBC-LP+F}  & $33.73$ & $0.107$  & $3586.6$\\
LSS                & $37.68$ & $18.254$ & $2211.5$\\
\texttt{CATZ-F4}   & $92.83$ & $74.304$ & $3506.5$\\
\bottomrule
\end{tabular}
\caption{Solving times on the $535$ instances that all eight configurations solve. The
Benders solver has by far the lowest median and a high mean, which is the signature of
a method that either finishes almost at once or struggles.}
\label{tab:times}
\end{table}

The pattern in Table~\ref{tab:times} is as informative as the ranking. The Benders
solver's median is one to three orders of magnitude below every compact model's, while
its mean is an order of magnitude above SSPMF's. It either closes an instance
immediately or does not close it at all. Section~\ref{ssec:results-tight} identifies
what separates the two cases, and it is not the size of the instance.

\subsection{The Benders solver and its strengthenings}\label{ssec:results-bbc}

\paragraph{Fractional cuts.}
Section~\ref{ssec:accel} recorded that fractional separation had been inert in the
earlier campaign because of a defect in the cut-injection call, so the idea could not
be tested. The defect is corrected, and the corrected version was run.

The cuts are now generated in very large numbers. Across the $1{,}415$ instances that
both \texttt{BBC-LP} and \texttt{BBC-LP+F} attempted, the fractional configuration
added $67{,}513{,}396$ cuts, and the configuration without fractional separation added
none, which is the expected control.

They do not help. Table~\ref{tab:frac} gives the comparison.

\begin{table}[h]\centering
\footnotesize
\begin{tabular}{@{}lrr@{}}
\toprule
On the $1{,}415$ common instances & \texttt{BBC-LP} & \texttt{BBC-LP+F}\\
\midrule
fractional cuts added        & $0$   & $67{,}513{,}396$\\
solved to optimality         & $726$ & $659$\\
median nodes explored        & $77{,}504$ & $3{,}795$\\
\midrule
\multicolumn{3}{@{}l}{final dual bound, on the $1{,}404$ instances where both are recorded:}\\
\quad strictly higher with fractional cuts & \multicolumn{2}{r}{$0$}\\
\quad equal                                & \multicolumn{2}{r}{$1{,}228$}\\
\quad strictly lower with fractional cuts  & \multicolumn{2}{r}{$176$}\\
\bottomrule
\end{tabular}
\caption{The effect of enabling fractional Benders cuts. Sixty-seven million cuts
enter the master. The bound at termination is never higher than without them, and on
$176$ instances it is lower, because the hour is spent generating and managing cuts
instead of searching. Sixty-seven instances that the base solver closes are lost.}
\label{tab:frac}
\end{table}

This is a clean negative result and it is worth stating plainly. The mechanism works:
the cuts are valid, they are generated, and they enter the master. The node counts show
they do tighten the relaxation locally, since the search explores twenty times fewer
nodes. What they do not do is raise the bound that matters. On no instance in the
comparison is the bound at termination higher with them than without.

The reading is that the fractional dual information is not new information. The
integer cuts already give the master everything the loading dual has to say, and the
fractional ones re-derive it at a cost of an hour of separation.

\paragraph{The three strengthenings.}
Table~\ref{tab:ablation} compares each strengthening against \texttt{BBC-LP+F} on the
instances both attempted.

\begin{table}[h]\centering
\footnotesize
\begin{tabular}{@{}llrrr@{}}
\toprule
Configuration & what it adds & common & solved & difference\\
\midrule
\texttt{BBC-LP+F}   & ---                        & ---       & ---   & ---\\
\texttt{BBC-LP+F+H} & hybrid-genetic incumbent   & $1{,}421$ & $690$ & $+31$\\
\texttt{BBC-LP+F+C} & conflict-graph row         & $1{,}421$ & $657$ & $-2$\\
\texttt{BBC-LP+ACC} & all three together         & $1{,}421$ & $656$ & $-3$\\
\texttt{BBC-LP+F+P} & Pareto-optimal cut lifting & $1{,}421$ & $622$ & $-37$\\
\midrule
\texttt{BBC-LP}     & fractional separation removed & $1{,}415$ & $726$ & $+67$\\
\bottomrule
\end{tabular}
\caption{Each strengthening against \texttt{BBC-LP+F}, on the instances both attempted.
Only the one that improves the incumbent helps. The two that aim at the dual bound
cost more time than they return, and combining all three does not recover the loss.
The last row is the comparison of Table~\ref{tab:frac}, repeated here because removing
a feature outperforms every addition.}
\label{tab:ablation}
\end{table}

The pattern is consistent and it is the same pattern as the fractional cuts. The only
strengthening that helps is the one that supplies a better \emph{incumbent}. Every
attempt to improve the \emph{bound} costs more than it returns. The conflict-graph row
is neutral, which is what Section~\ref{ssec:accel} predicted for it on structural
grounds: on most instances the coverage row already dominates it.

\paragraph{Where the search stops.}
Two further measurements locate the difficulty precisely.

First, of the $726$ instances \texttt{BBC-LP} solves, $259$ are closed at the root
node, without branching at all. The root relaxation value was recorded on $1{,}178$ of
the $1{,}421$ runs, and on $883$ of those the optimum is also known. On those $883$ the
root value equals the optimum $294$ times, and the split is sharp: among the runs that
went on to close the instance it is already the optimum in $264$ of $497$, and among
the runs that did not close it in only $30$ of $386$. Where it is not the optimum it is
short of it by at least one unit and by $16.7\%$ at the median, rising to $42.9\%$.

Second, consider the runs that hit the time limit. Across the nine Benders
configurations there are $6{,}562$ of them, and for $3{,}874$ the optimum is known
because some method proved it. On $2{,}571$ of those $3{,}874$, which is $66\%$, the
Benders solver's incumbent already \emph{equals} that optimum. It had found the answer
and could not prove it. For \texttt{BBC-LP} alone the figure is $229$ of $386$, or
$59\%$.

Taken together: when the bound is good the solver finishes at the root, and when the
bound is bad it finds the optimum quickly and then spends an hour failing to certify
it. The search is not the bottleneck. The bound is.

\subsection{The coverage bound decides the outcome}\label{ssec:results-tight}

Section~\ref{sec:gap} predicted that the coverage bound $q$ would separate easy
instances from hard ones for any configuration-based method. That prediction is
directly testable, because for each instance with a known optimum it can be checked
whether the optimum equals the bound.

Of the $1{,}115$ instances whose optimum is known, $523$ have the coverage bound tight
and $592$ do not. The split is close to even, at $46.9\%$ tight, which means the
standard instance collections are not concentrated in either regime. Where the bound is
loose it is loose by a wide margin: the optimum exceeds it by $5.71$ on average and by
as much as $30$.

Splitting the results by that criterion gives Table~\ref{tab:tightness}.

\begin{table}[h]\centering
\footnotesize
\begin{tabular}{@{}lrrrr@{}}
\toprule
& instances & \texttt{BBC-LP} solves & SSPMF solves\\
\midrule
coverage bound tight ($\Zst=q$)  & $523$ & $492$ \; ($94\%$) & $523$ \; ($100\%$)\\
coverage bound loose ($\Zst>q$)  & $592$ & $234$ \; ($40\%$) & $515$ \; ($87\%$)\\
\bottomrule
\end{tabular}
\caption{Success rate split by whether the coverage bound equals the optimum. The
bound, not the size of the instance, is what decides the outcome for the Benders
solver. The multicommodity model is affected by the same split but far less, because
it is not relying on that bound alone.}
\label{tab:tightness}
\end{table}

This is the central measurement of the campaign. The same solver, on the same hardware,
under the same time limit, closes $94\%$ of one class and $40\%$ of the other. The
classification is not by number of jobs, number of tools or density; it is by a single
structural property of the instance, and it is the property the theory of
Section~\ref{sec:gap} identifies.

The SSPMF column is the useful control. Its relaxation is also equal to $q$, so if the
bound were the whole story it should show the same collapse. It does not: it falls from
$100\%$ to $87\%$ where the Benders solver falls from $94\%$ to $40\%$. A compact model
gives the branch-and-bound search a full description of the problem to work with at
every node, and the search recovers much of what the bound does not supply. The Benders
master sees only the cuts it has accumulated, and when the bound is loose it has very
little else.

\subsection{Branch-and-price}\label{ssec:results-bnp}

The two position-indexed prototypes were run separately from the campaign above. They
are implemented in Python over PySCIPOpt with a Python pricer, so a single node costs
orders of magnitude more than a node of a compiled solver, and comparing their
wall-clock times against CPLEX would measure the implementation language rather than
the formulation. In total $2{,}581$ runs were made over $469$ instances in seven
configurations: PCF$'$ alone; PCF$'$ with each of the four accelerations of
Section~\ref{ssec:pos} separately, written \texttt{+HP} for heuristic pricing,
\texttt{+MC} for multiple pricing, \texttt{+WS} for warm starting and \texttt{+STAB}
for dual stabilisation; PCF$'$ with all four together, written \texttt{+ACC}; and PTF.
The time limit is thirty minutes on the Catanzaro and Crama families and ten minutes on
the Laporte families.

Two properties of this run make its counts weaker evidence than the campaign above, and
both are stated before the numbers. First, the harness carries an early-stop rule: once
a configuration has failed to close six increasingly hard instances in succession, the
remaining harder instances are skipped for that configuration. The number of runs a
configuration received therefore depends on how it performed, so the denominators in
Table~\ref{tab:bnp} are not comparable across rows. Second, no repetition or seed
variation was run. Every comparison drawn below is made either on instances that both
configurations attempted, or on the root relaxation value, which the early-stop rule
cannot bias because it is recorded before any search.

\paragraph{They agree with the campaign.}
On the $507$ runs that proved an optimum and where the same instance was also closed in
the campaign of Section~\ref{ssec:results-solved}, the two agree on every one. As in
Section~\ref{ssec:validation}, the native objective is the free-initial cost and the
difference from the empty-start cost is exactly $b$ on all $507$.

\paragraph{What they solve.}
Table~\ref{tab:bnp} gives the counts. Nothing above nine jobs was closed by any
configuration: all $608$ runs on the $30$-job and $40$-job instances reached the time
limit. This is the expected consequence of a Python pricer and is not evidence about
the formulation.

\begin{table}[h]\centering
\footnotesize
\begin{tabular}{@{}llrrrr@{}}
\toprule
Configuration & what it adds & runs & solved & median time (s) & median columns\\
\midrule
PCF$'$\texttt{+MC}   & multiple pricing          & $422$ & $102$ & $11.4$ & $1{,}504$\\
PCF$'$               & ---                       & $415$ & $91$  & $23.7$ & $1{,}574$\\
PCF$'$\texttt{+HP}   & heuristic pricing         & $360$ & $73$  & $19.7$ & $5{,}173$\\
PTF                  & position--transition rows & $387$ & $73$  & $2.4$  & $1{,}279$\\
PCF$'$\texttt{+WS}   & warm-started pool         & $359$ & $64$  & $14.9$ & $1{,}610$\\
PCF$'$\texttt{+ACC}  & all four together         & $337$ & $60$  & $28.1$ & $1{,}221$\\
PCF$'$\texttt{+STAB} & dual stabilisation        & $301$ & $44$  & $35.7$ & $1{,}133$\\
\bottomrule
\end{tabular}
\caption{The branch-and-price prototypes. Median time is over the runs that closed;
median columns is over all runs. The run counts are outcome-dependent through the
early-stop rule, so the rows are not comparable; the comparison that matters is the
one below, which is made on common instances.}
\label{tab:bnp}
\end{table}

\paragraph{The measurement that matters.}
Proposition~\ref{prop:pcf} states that the position--configuration relaxation equals
the coverage bound $q$, and Proposition~\ref{prop:ptf} states that the
position--transition relaxation can exceed it. Both are testable directly, because the
prototypes record the root relaxation value.

\emph{PCF$'$ never exceeds $q$.} On all $2{,}114$ PCF$'$ runs for which a root
relaxation value was recorded, that value equals $\lvert\U\rvert-b$ exactly. Not one
run, in any of the six PCF$'$ configurations, produced a root bound above the counting
bound. This is Proposition~\ref{prop:pcf} confirmed on the benchmark rather than on
paper.

\emph{PTF exceeds $q$, and almost never.} Of the $233$ PTF runs with a recorded root
relaxation, three are strictly above $\lvert\U\rvert-b$: Catanzaro \texttt{A0-0}, where
the bound is $2.013$ against $q=2$, and Laporte3 \texttt{L11-6} and \texttt{L11-7},
where it is $6$ against $q=5$. Proposition~\ref{prop:ptf} is therefore not vacuous. But
$230$ of $233$ runs sit exactly on $q$, and on the two Laporte3 instances the optima
are $12$ and $13$, so a bound of $6$ instead of $5$ closes a small part of a large gap.

\emph{The bound is often already the answer.} On $990$ of the $1{,}846$ runs where both
the root relaxation and the optimum are known, the root relaxation \emph{equals} the
optimum. These are the tight instances of Section~\ref{ssec:results-tight}, seen from
the other side: where the coverage bound is tight, a configuration-based relaxation
needs no strengthening at all, and where it is loose, neither PCF$'$ nor PTF supplies
one.

Head to head on the $320$ instances both attempted, PCF$'$ and PTF each close $73$.
Their node counts differ by a factor of fifty --- a median of $96$ against $2$ ---
because the PTF pricer prices coupled pairs of configurations and costs far more per
node. The stronger relaxation does not convert into more instances closed.

\paragraph{The comparison with the Job Grouping Problem.}
\citet{Otiai2024} builds a branch-and-price for that problem whose set-covering
relaxation is strong enough to close benchmark instances in a handful of nodes. For the
SSP, \citet{Moreira2026} reports that the trivial bound $\lvert T\rvert-b$ obtained by
the multicommodity model was the best available bound on all thirty instances of three
of his subgroups, with the configuration model giving the best bound on four instances
of the remaining, tighter subgroup. The measurements above say the same thing on a
larger collection and with the mechanism visible: the configuration relaxation for the
SSP is not, in general, an improvement on the counting bound.

\subsection{What the campaign establishes}\label{ssec:results-summary}

Six statements are supported by the measurements above.

\begin{enumerate}[topsep=4pt,itemsep=3pt]
  \item The reimplemented multicommodity model of \citet{daSilva2024} is the strongest
  single method in this study, which replicates the claim of its authors on an
  independent implementation and a different instance collection.

  \item Cross-method agreement is complete. On the $999$ instances solved by more than
  one method, all $41{,}673$ pairs of methods agree on the cost of the schedule
  returned.

  \item Fractional Benders cuts, once the implementation defect is corrected, are
  generated in very large numbers and are counterproductive. Sixty-seven million of
  them never raise the bound at termination on a single instance, and they cost
  sixty-seven solved instances.

  \item Of four strengthenings of the Benders solver, only the one that improves the
  incumbent helps. The three that aim at the bound are neutral or harmful.

  \item Whether the coverage bound equals the optimum decides whether the Benders
  solver closes the instance: $94\%$ against $40\%$. Where it fails, the incumbent is
  already optimal $59\%$ of the time.

  \item The position--configuration relaxation equals the coverage bound on every one
  of $2{,}114$ runs, and the position--transition relaxation exceeds it on three runs of
  $233$. The first is Proposition~\ref{prop:pcf} confirmed; the second is
  Proposition~\ref{prop:ptf} confirmed and, at the same time, shown to be of little
  practical use as it stands.
\end{enumerate}

\noindent
The first two are checks that the study is sound. The last four all say the same thing,
which Section~\ref{sec:disc} takes up.
